\ifdefined\iscombined
	\lecturetitle{Mechanical Systems}
\else
\documentclass{article}
\usepackage{xparse}
\usepackage{changepage}
\usepackage{listings}
\usepackage{amsthm}
\usepackage{comment}
\usepackage{amsmath}
\usepackage{braket}
\usepackage{amsfonts}
\usepackage{sectsty}
\usepackage{hyperref}
\usepackage{fancyhdr}
\usepackage[top=1in,bottom=1in,left=1in,right=1in]{geometry}
\usepackage{float}
\usepackage{graphicx}
\usepackage{mathtools}
\usepackage{tikz}
\usetikzlibrary{arrows}
\usepackage{circuitikz}

\date{
	\vspace{-.75em}
	{\large \today}
}

\title{
	\vspace*{-1.5em}
	{\Huge Lecture 9} \\
	\vspace{-1em}
	\rule{\textwidth/4}{.01em} \\
	\vspace{-.25em}
	{\Large Mechanical Systems}
	\vspace{-.75em}
}

\author{
	{\large Matthew Farah}
}

\hypersetup{
	colorlinks=true,
	linkcolor=blue,
	filecolor=magenta,
	urlcolor=blue,
	citecolor=blue,
	anchorcolor=blue
}

\fancyhead{}
\fancyhead[L]{\today}
\fancyhead[R]{Lecture 9 - Mechanical Systems}

\setlength{\parindent}{0cm}

\newcounter{example}
\renewcommand{\theexample}{\arabic{example}}

\newenvironment{example}[1][]{
	\refstepcounter{example}
	\begin{adjustwidth}{1.5em}{}
	\noindent\textbf{\ifx&#1&Example \theexample:\else Example \theexample \ (#1):\fi}
	\ignorespaces
}{
	\vspace{.5em}
	\end{adjustwidth}
}

\newenvironment{solution}{
	\vspace{.5em}
	\par
	\begin{adjustwidth}{1.5em}{}
	\textbf{{Solution:}}
	\par
	\vspace{.5em}
}{
	\vspace{1em}
	\end{adjustwidth}
}

\newcounter{subexample}
\renewcommand{\thesubexample}{\alph{subexample}}

\newenvironment{subexample}{
	\setcounter{step}{0}
	\refstepcounter{subexample}
	\vspace{1em}\par\noindent
	\begin{adjustwidth}{1.5em}{}
	\noindent\textbf{(\thesubexample)}
	\ignorespaces
}{
	\par
	\vspace{1em}
	\end{adjustwidth}
}

\renewenvironment{proof}{
	\vspace{.5em}\par\noindent
	\begin{adjustwidth}{1.5em}{}
	\emph{Proof:}\par\vspace{1em}
	\setlength{\parindent}{0cm}
}{
	\end{adjustwidth}
}

\newtheorem{theorem}{Theorem}

\renewenvironment{theorem}[1][]{
	\refstepcounter{theorem}
	\vspace{1em}\par\noindent
	\begin{adjustwidth}{1.5em}{}
	\noindent\textbf{\ifx&#1&Theorem \thetheorem:\else Theorem \thetheorem \ (#1):\fi}
	\ignorespaces
}{
	\vspace{.5em}
	\end{adjustwidth}
}

\newtheorem{lemma}{Lemma}

\renewenvironment{lemma}[1][]{
	\refstepcounter{lemma}
	\vspace{1em}\par\noindent
	\begin{adjustwidth}{1.5em}{}
	\noindent\textbf{\ifx&#1&Lemma \thelemma:\else Lemma \thelemma \ (#1):\fi}
	\ignorespaces
}{
	\vspace{.5em}
	\end{adjustwidth}
}

\makeatother
\makeatletter

\newcounter{step}
\renewcommand{\thestep}{\Roman{step}}

\newenvironment{step}{
	\refstepcounter{step}
	\vspace{1em}\par\noindent
	\ignorespaces\noindent\textbf{(\thestep)}
	\ignorespaces
}{
	\par
}

\sectionfont{\fontsize{12}{12}\bfseries}
\subsectionfont{\fontsize{10}{12}\normalfont}
\subsubsectionfont{\fontsize{10}{12}\normalfont}

\begin{document}

\maketitle
\pagestyle{fancy}

\fi

\begin{itemize}
	\item The two kinds of mechanical systems we'll analyze in this course are:
	\begin{enumerate}
		\item Translational systems.
		\item Rotational systems.
	\end{enumerate}
	\item The transfer function of a mechanical system can be implemented as a circuit and vice-versa.
	\item The transfer functions of translational systems describe the ratio between an input force and the displacement of some body.
	\item Impedances in translational systems work as follows:
	\begin{itemize}
		\item A mass $M$ imparts an impedance of $Ms^2$ on a moving body.
		\item A viscous damper with damping constant $f_v$ imparts an impedance of $f_v s$ on a moving body.
		\item A spring with spring constant $K$ imparts an impedance of $K$ on a moving body.
	\end{itemize}
	\item The direction of a body's motion is defined relative to the external force acting on the system.
	\item Newton's third law of motion states that the net force acting on an object must sum to zero. In other words: \\
\begin{align*}
	\sum_{F_n \in \text{object}} F_n &= 0 \\
\end{align*}
\end{itemize}

\begin{example}
	Draw the Free Body Diagram (FBD) and determine the transfer function $G(s)$ for a system consisting of a body:
	\begin{itemize}
		\item With a mass $M$.
		\item Attached to a viscous damper with damping coefficient $f_v$.
		\item Attached to a spring with spring constant $K$.
		\item Being accelerated by a force $f(t)$.
	\end{itemize}
\end{example}

\begin{solution}
	Based on the description given, we can depict the system in the time domain like so:

	\vspace{1em}

	\begin{figure}[H]
		\centering
		\begin{tikzpicture}[>=latex']
			\draw (0,-1.4) rectangle (1.5,2.8);
			\node at (0.75,0.7) {$M$};
			\draw[->] (0,2.4) -- node[midway,above]{$M \ddot{x}(t)$} (-3.2,2.4);
			\draw[->] (0,1.4) -- node[midway,above]{$Kx(t)$} (-3.2,1.4);
			\draw[->] (0,0.4) -- node[midway,above]{$f_v \dot{x}(t)$} (-3.2,0.4);
			\draw[->] (1.5,1.4) -- node[midway,above]{$f(t)$} (4.7,1.4);
			\draw[dashed, ->] (0.75,3.3) -- node[midway,above]{$x(t)$} (2.75,3.3);
		\end{tikzpicture}
	\end{figure}

	\vspace{1em}

	In the Laplace domain, this looks like:

	\vspace{1em}

	\begin{figure}[H]
		\centering
		\begin{tikzpicture}[>=latex']
			\draw (0,-1.4) rectangle (1.5,2.8);
			\node at (0.75,0.7) {$M$};
			\draw[->] (0,2.4) -- node[midway,above]{$Ms^2 X(s)$} (-3.2,2.4);
			\draw[->] (0,1.4) -- node[midway,above]{$KX(s)$} (-3.2,1.4);
			\draw[->] (0,0.4) -- node[midway,above]{$f_v s X(s)$} (-3.2,0.4);
			\draw[->] (1.5,1.4) -- node[midway,above]{$F(s)$} (4.7,1.4);
			\draw[dashed, ->] (0.75,3.3) -- node[midway,above]{$X(s)$} (2.75,3.3);
		\end{tikzpicture}
	\end{figure}

	\vspace{1em}

	We know the transfer function for the body's displacement, $G(s)$, is given by $G(s) = \frac{X(s)}{F(s)}$. By Newton's third law of motion, the FBD above implies that:

	\begin{equation*}
		\begin{array}{lrl}
			& KX(s) + f_v s X(s) + Ms^2 X(s) =& F(s) \\
			\implies & (K + f_v s + Ms^2)X(s) =& F(s) \\
			\implies & X(s) =& \dfrac{F(s)}{Ms^2 + f_v s + K} \\\\
		\end{array}
	\end{equation*} \\

	And thus:

	\begin{align*}
		G(s) &= \frac{X(s)}{F(s)} = \frac{1}{Ms^2 + f_v s + K}
	\end{align*}
\end{solution}

\begin{itemize}
	\item The degrees of freedom (dof) of a translational system describes the number of linearly independent motions within the system.
	\begin{itemize}
		\item Two bodies have linearly independent motions from each other if you're still able to move one of them if you keep the other locked in place.
		\item The number of equations needed to find the transfer function of a system is equal to its degrees of freedom.
	\end{itemize}
	\item Mechanical systems obey the superposition principle.
	\begin{itemize}
		\item States that the equation describing the motion of a body $M_i$ can be found by:
		\begin{enumerate}
			\item Holding all other bodies still and finding the forces acting on $M_i$ solely due to its own motion.
			\item Holding $M_i$ and all other bodies still except for some body $M_j$, and finding all forces acting on $M_i$ induced by moving $M_j$.
			\item Repeating 2) until every body with a linearly independent motion to $M_i$ has been covered.
			\item Summing all forces acting on $M_i$ found in steps 1) to 3) and setting that equal to zero.
		\end{enumerate}
	\end{itemize}
\end{itemize}

\ifdefined\iscombined
	\newpage
\else
	\end{document}
\fi
